Safety evaluations and mechanistic analyses of refusal usually assume that a model's response to a harmful
request begins at its first output token. We show that this assumption breaks for prompts that bundle a harmful
request with benign ones, and that it produced a false result in our own earlier work. There, scoring the first
24 generated tokens suggested that a numbered list of benign tasks collapses Qwen3-14B's refusal after about 31
tokens, and probes at the last prompt token appeared to explain why. We re-measure this setting with the harmful
item's position controlled, full-length generations, and whole-response judging: three model families, 200
JailbreakBench and AdvBench targets, lists of up to 32 items, three templates, about 18{,}000 responses. When the
harmful item is not first, the 24-token judge reports refusal rates as low as \QwenMiddleVone\ (Qwen3-14B) and
\GemmaMiddleVone\ (Gemma-2-9B). Yet the models refuse the item in \QwenMiddleRef\ and \GemmaMiddleRef\ of those
responses, and HarmBench's classifier finds harmful content in at most 15\% of responses in any
condition. Models answer the benign items and refuse when they reach the harmful one (Qwen, Gemma), or refuse it
up front (Llama). At the last prompt token, list structure shrinks the refusal-direction projection and can flip its sign
(16\% of Qwen3-14B list prompts). Yet the projection still ranks responses by whether the model will eventually refuse the item
(AUC up to \QwenAUCLastBest), and the same information is readable at the item's answer onset. Adding the
unnormalized refusal direction at that point lowers harmful compliance. Short-horizon judges and last-token
probes are therefore unreliable for multi-request prompts, and we give a protocol that avoids both problems.
